\section{Decision Episode Model}

\subsection{Purpose}

A Decision Episode represents a significant engineering decision together with the context that produced it.

Its purpose is to preserve engineering reasoning rather than merely recording the final decision.

Unlike traditional documentation, a Decision Episode captures how a decision evolved, why it was made, what alternatives were considered, and what evidence supported the final outcome.

\vspace{0.5cm}

\subsection{Motivation}

Engineering projects involve numerous design decisions throughout their lifecycle.

Although source code, documentation, and version control preserve the final implementation, they rarely preserve the reasoning that produced it.

Consequently, engineers frequently spend considerable effort reconstructing historical context during design reviews, debugging sessions, onboarding, or discussions with collaborators.

The Decision Episode Model addresses this limitation by treating engineering decisions as first-class project artifacts.

\vspace{0.5cm}

\subsection{Definition}

A Decision Episode is an interpreted representation of one or more related engineering activities that together describe a meaningful engineering decision.

A Decision Episode is not created from a single event.

Instead, it is constructed by correlating evidence collected throughout the engineering workflow.

\vspace{0.5cm}

\subsection{Decision Episode Lifecycle}

A Decision Episode progresses through the following stages.

\begin{enumerate}
    \item Evidence Collection
    \item Evidence Correlation
    \item Candidate Decision Episode
    \item Confidence Assessment
    \item Engineer Clarification (when required)
    \item Confirmed Decision Episode
    \item Project Understanding Update
\end{enumerate}

\vspace{0.5cm}

\subsection{Evidence Sources}

Evidence supporting a Decision Episode may originate from multiple sources.

Typical sources include:

\begin{itemize}
    \item Engineer input
    \item Source code evolution
    \item Version control history
    \item RTL architecture modifications
    \item Testbench evolution
    \item Simulation results
    \item Synthesis reports
    \item Guide feedback
    \item Research papers
    \item Engineering notes
    \item Design documentation
\end{itemize}

The system is not required to utilize every available source.

The objective is to accumulate sufficient evidence to construct a useful understanding of the engineering decision.

\vspace{0.5cm}

\subsection{Decision Episode Components}

Each confirmed Decision Episode should maintain the following information.

\begin{itemize}
    \item Trigger
    \item Engineering Problem
    \item Project Context
    \item Alternatives Considered
    \item Selected Solution
    \item Engineering Reasoning
    \item Supporting Evidence
    \item Expected Outcome
    \item Actual Outcome (when available)
    \item Remaining Open Questions
\end{itemize}

\vspace{0.5cm}

\subsection{Confidence Assessment}

Decision Episodes are created from interpreted evidence and therefore possess varying levels of confidence.

The system shall estimate its confidence before accepting a candidate Decision Episode.

Three confidence levels are defined.

\begin{itemize}
    \item High Confidence
    \item Medium Confidence
    \item Low Confidence
\end{itemize}

High-confidence episodes may be accepted automatically.

Medium and low-confidence episodes should request clarification from the engineer before becoming part of the Project Understanding Model.

\vspace{0.5cm}

\subsection{Engineer Interaction}

The objective of the Decision Episode Model is to minimize interruption while preserving engineering understanding.

The engineer should not be required to manually document every decision.

Instead, the Shadow should:

\begin{enumerate}
    \item Observe engineering activities.
    \item Correlate available evidence.
    \item Construct a candidate Decision Episode.
    \item Request clarification only when necessary.
    \item Update the Project Understanding Model.
\end{enumerate}

Consequently, engineer interaction should remain brief and focused on validating or correcting inferred understanding rather than writing documentation.

\vspace{0.5cm}

\subsection{Relationship with the Project Understanding Model}

Decision Episodes do not exist independently.

Every confirmed Decision Episode contributes to the continuously evolving Project Understanding Model.

The Project Understanding Model represents the Shadow's current understanding of the project, while Decision Episodes explain how that understanding evolved over time.

\vspace{0.5cm}

\subsection{Design Principles}

The Decision Episode Model follows the principles below.

\begin{enumerate}
    \item Preserve reasoning rather than only outcomes.
    \item Correlate evidence from multiple sources.
    \item Minimize engineer interruption.
    \item Prefer deterministic interpretation whenever sufficient.
    \item Utilize AI reasoning only when deterministic interpretation cannot provide sufficient understanding.
    \item Maintain transparency regarding confidence and supporting evidence.
    \item Continuously refine project understanding as new evidence becomes available.
\end{enumerate}

\vspace{0.5cm}

\subsection{Current Limitations}

The Decision Episode Model does not attempt to reconstruct the complete cognitive state of the engineer.

Instead, it aims to maintain a practical level of project understanding that provides meaningful engineering assistance while remaining feasible with current technologies.

The objective is useful project understanding rather than perfect human-level reasoning.